\cgeentry{cge_shock}{Change a parameter or a fixed variable}

\cgesec{Description}

\cmd{cge_shock} changes a parameter of the calibrated model, or a variable that
the current closure fixes, before the next \manref{cge_solve}. Shocks
accumulate until \cmd{cge_shock, reset} restores every base value.

\cgesec{Quick start}

Remove every import tariff
\begin{stcode}
cge_shock tm, value(0)
\end{stcode}
Raise the world import price of manufactures by 10 per cent
\begin{stcode}
cge_shock pwm, elements(com_man) scale(1.1)
\end{stcode}
Raise government purchases of health services by 20 per cent
\begin{stcode}
cge_shock qg, elements(com_hlt) scale(1.2)
\end{stcode}
Raise the supply of labour, which the default closure fixes, by 5 per cent
\begin{stcode}
cge_shock QFS, elements(lab) scale(1.05)
\end{stcode}
Remove every shock
\begin{stcode}
cge_shock, reset
\end{stcode}

\cgesec{Syntax}

Change a parameter or a fixed variable
\begin{stsyntax}
\cmdb{cge_shock} \argu{name} [\texttt{,} \optarg{el}{ements}{list}]
\texttt{\{}\optarg{sc}{ale}{\#}\,|\,\optarg{val}{ue}{\#}\texttt{\}}
\end{stsyntax}
Remove every shock
\begin{stsyntax}
\cmdb{cge_shock}\texttt{,} \opt{}{reset}
\end{stsyntax}
\argu{name} is a parameter, in lower case, from the table under
\emph{Remarks and examples}, or a model variable, in upper case, that the
current closure fixes.

\begin{opttable}
\optgroup{Main}
\optarg{el}{ements}{list} & elements to change; default all\\
\optarg{sc}{ale}{\#} & multiply the current values by \argu{\#}\\
\optarg{val}{ue}{\#} & set the values to \argu{\#}\\
\opt{}{reset} & remove every shock\\
\end{opttable}

Exactly one of \opt{sc}{ale()} and \opt{val}{ue()} is required, except with
\opt{}{reset}, which may not be combined with a shock.

\cgesec{Options}

\begin{optdesc}
\item[\optarg{el}{ements}{list}] lists the elements to change: commodities,
activities, factors, or households and enterprises, as the parameter or
variable requires; for a variable added by \manref{cge_variable}, the
elements of the set of its index, such as institutions for an index \cmd{i}.
The default is every element, as is \cmd{elements(all)}. For
a variable indexed by factor and activity, $\V{QF}_{fa}$ or
$\V{WFDIST}_{fa}$, the list names factors, and every activity of each factor is
changed. A parameter or variable without index takes no elements.

\item[\optarg{sc}{ale}{\#}] multiplies the current values by \argu{\#}. The
current values include earlier shocks, so two shocks of \cmd{scale(1.1)} give
a rise of 21 per cent.

\item[\optarg{val}{ue}{\#}] sets the values to \argu{\#}.

\item[\opt{}{reset}] removes every shock: all parameters and fixed variables
return to their base values. The constraints of \manref{cge_constraint}, and
the recipients of their rents, are kept.
\end{optdesc}

\cgesec{Remarks and examples}

A simulation compares an equilibrium under changed parameters or fixed values
with the base. \cmd{cge_shock} makes the changes; \manref{cge_solve} then
solves the model under them and stores the simulation, with the parameters it
used. The shocks stay in memory after the solve, so that further shocks add to
them, until \cmd{cge_shock, reset}. \manref{cge_calibrate} also removes them.

\cgesubsec{Parameters}

The parameters that \cmd{cge_shock} changes are listed below with the set over
which each is defined; those marked ``no elements'' are scalars. Values in
foreign currency are converted at the exchange rate, and values ``in real
terms'' are indexed to the CPI. The symbols are those of
\secref{sec:parameters}.

\begin{cgetable}{p{0.08\linewidth}p{0.26\linewidth}p{0.59\linewidth}}
\toprule
Name & Elements & Description\\
\midrule\endhead
\bottomrule\endlastfoot
\multicolumn{3}{@{}l}{\textit{Taxes}}\\
\cmd{tm} & commodities & import tariff rates, $\V{tm}_c$\\
\cmd{te} & commodities & export tax rates, $\V{te}_c$\\
\cmd{tq} & commodities & sales tax rates, $\V{tq}_c$\\
\cmd{ta} & activities & activity tax rates, $\V{ta}_a$\\
\cmd{tv} & commodities & VAT rates, $\V{tv}_c$; needs a VAT account\\
\cmd{kappa} & activities & VAT credit shares, $\kappa_a$, between 0 and 1; needs a VAT account\\
\cmd{tf} & factors & rates of tax on the use of factors, $\V{tf}_{fa}$, in every activity that uses them\\
\cmd{tfy} & factors & rates of tax on factor income, $\V{tfy}_f$\\[2pt]
\multicolumn{3}{@{}l}{\textit{World prices}}\\
\cmd{pwm} & commodities & world import prices, $\V{pwm}_c$, in foreign currency\\
\cmd{pwe} & commodities & world export prices, $\V{pwe}_c$; with finite export demand, a shift in world demand\\[2pt]
\multicolumn{3}{@{}l}{\textit{Productivity}}\\
\cmd{ad} & activities & efficiency of value added, $\V{ad}_a$, that is productivity\\
\cmd{af} & factors & efficiency of each factor in every activity, $\V{af}_f$; one in the base\\[2pt]
\multicolumn{3}{@{}l}{\textit{Spending and transfers}}\\
\cmd{qg} & commodities & base quantities of government consumption, $\overline{\V{qg}}_c$, which $\V{GADJ}$ scales\\
\cmd{qinv} & commodities & base quantities of investment, $\overline{\V{qinv}}_c$, which $\V{IADJ}$ scales\\
\cmd{trdg} & households and enterprises & government transfers, $\V{trnsfr}_{d,gov}$, in real terms\\
\cmd{trdr} & households and enterprises & transfers from the rest of the world, $\V{trnsfr}_{d,row}$, in foreign currency\\
\cmd{trfr} & factors & factor income from abroad, $\V{trnsfr}_{f,row}$, in foreign currency\\
\cmd{trgr} & no elements & transfers from the rest of the world to the government, $\V{trnsfr}_{gov,row}$, for example budget support, in foreign currency\\
\cmd{trrg} & no elements & government transfers to the rest of the world, $\V{trnsfr}_{row,gov}$, in foreign currency\\[2pt]
\multicolumn{3}{@{}l}{\textit{Health insurance fund (needs a fund, mapped with \cmd{cge_map, insurance()})}}\\
\cmd{tshi} & factors & contribution rates on factor income, $\V{tshi}_f$\\
\cmd{qshi} & commodities & base purchases of the fund, $\overline{\V{qshi}}_c$, which $\V{SHIADJ}$ scales\\
\cmd{trshi} & no elements & government transfers to the fund, $\V{trnsfr}_{shi,gov}$, in real terms\\
\cmd{trshir} & no elements & transfers from the rest of the world to the fund, $\V{trnsfr}_{shi,row}$, in foreign currency\\[2pt]
\multicolumn{3}{@{}l}{\textit{Public debt (needs debt, from \cmd{cge_calibrate, debt() interest()})}}\\
\cmd{rd}, \cmd{rf} & no elements & interest rates on domestic and foreign public debt, $r^d$ and $r^f$\\
\cmd{ddebt} & no elements & domestic public debt at the start of the period, $\V{DD}$, in local currency\\
\cmd{fdebt} & no elements & foreign public debt at the start of the period, $\V{DF}$, in foreign currency; for example debt relief\\[2pt]
\multicolumn{3}{@{}l}{\textit{Labour markets (factors with a labour force, from \cmd{cge_calibrate, unemployment()} or \cmd{minwage()})}}\\
\cmd{qlf} & factors & labour force, $\V{QLF}_f$\\
\cmd{wmin} & factors & real minimum wage, $\underline{w}_f$\\
\cmd{wref} & factors & reference real wage of the wage curve, $\overline{w}_f$\\[2pt]
\multicolumn{3}{@{}l}{\textit{Health outcomes (from \cmd{cge_calibrate, sickdays()} or \cmd{mortality()})}}\\
\cmd{sick} & factors & share of working time lost to illness before the response to health services, $\sigma_f$; for example an epidemic\\
\cmd{mort} & no elements & crude death rate before the response to health services, $\mu^m$\\
\end{cgetable}

A shock to a parameter that the model lacks is an error: \cmd{tv} and
\cmd{kappa} need a VAT account, mapped with \cmd{cge_map, vattax()} or
\cmd{vatshare()} (\cmd{vatshare(0)} creates an empty one, so that VAT can be
introduced); \cmd{rd}, \cmd{rf}, \cmd{ddebt} and \cmd{fdebt} need public debt;
\cmd{qlf}, \cmd{wmin} and \cmd{wref} need factors with a labour force;
\cmd{sick} needs days lost to illness for the factors given, and \cmd{mort} a
death rate. The values must stay in range: \cmd{kappa} between 0 and 1,
\cmd{sick} and \cmd{mort} at least 0 and below 1, and \cmd{qlf}, \cmd{wmin} and
\cmd{wref}, when set with \opt{val}{ue()}, positive.

Some parameters call for a remark. \cmd{qg} and \cmd{qinv} are the base
quantities that $\V{GADJ}$ and $\V{IADJ}$ scale, so changing some of them
changes the composition of government consumption or investment. In the
default closure government savings adjust to a change in government spending
or transfers. \cmd{af} multiplies the efficiency of a factor: factor quantities
and wages stay in physical units, so with mobile factors a fall of 2 per cent
in labour efficiency has the same real effects as 2 per cent less labour. When
\manref{cge_calibrate} links efficiency to health services, more health
services also raise the efficiency variable $\V{EFF}_f$. For \cmd{tf},
\opt{sc}{ale()} multiplies each rate of the factors given, and
\opt{val}{ue()} sets one rate for every activity that uses them, which
introduces a tax where the SAM has none. \cmd{tshi}, \cmd{qshi}, \cmd{trshi}
and \cmd{trshir} change the revenue of the health insurance fund, which spends
it less a fixed share saved. Under the closures \cmd{wagecurve} and
\cmd{minwage} of \manref{cge_closure}, a larger labour force raises
unemployment unless the demand for labour grows with it. Doubling \cmd{sick}
doubles the share of working time lost at base health services per head. A
higher death rate changes no result within a period; in a dynamic run it slows
the growth of population and labour from the next period.

\cgesubsec{Fixed variables}

A variable can be changed while the current closure fixes it; changing a
variable that the closure leaves free is an error. Only the elements that
exist in the model are changed. Any variable element fixed by a swap of
\manref{cge_closure} can be changed, including a user variable defined by
\manref{cge_variable}. The closure options fix the following variables:

\begin{cgetable}{p{0.30\linewidth}p{0.66\linewidth}}
\toprule
Variable & Fixed under\\
\midrule\endhead
\bottomrule\endlastfoot
\cmd{FSAV}, \cmd{EXR} & \cmd{external(flexexr)}, \cmd{external(fixexr)}\\
\cmd{GSAV} & \cmd{government(fixgsav)}\\
\cmd{TINSADJ}, \cmd{TQADJ}, \cmd{TMADJ}, \cmd{TVADJ} & every government closure, except the factor that \cmd{fixgsav} frees\\
\cmd{GADJ} & every closure\\
\cmd{MPSADJ}, \cmd{IADJ} & \cmd{savinv(savings)}, \cmd{savinv(investment)}\\
\cmd{CPI}, \cmd{DPI} & \cmd{numeraire(cpi)}, \cmd{numeraire(dpi)}\\
\cmd{QFS} & factors that are \cmd{mobile}\\
\cmd{W} & factors that are \cmd{specific} or \cmd{unemployed}\\
\cmd{QF} & factors that are \cmd{specific}\\
\cmd{WFDIST} & factors that are \cmd{mobile}, \cmd{unemployed}, \cmd{wagecurve} or \cmd{minwage}\\
\cmd{RM}, \cmd{RE}, \cmd{RA} & every closure, unless a constraint frees them\\
\end{cgetable}

The tax scaling factors can therefore be set while the closure fixes them. A premium on imports or exports, or a rent on
output, set by hand acts as a tax whose revenue goes to the recipient of the
rent, the government unless \manref{cge_constraint} names another. A shock to
the adjusting variable of a constraint changes the level at which it is fixed
while the constraint does not bind. Of the variables with two indexes, only
$\V{QF}_{fa}$ and $\V{WFDIST}_{fa}$ can be shocked.

\cgesubsec{Example 1: parameters}

The synthetic SAM is mapped and calibrated as in \manref{cge_solve}. Each shock
reports the number of elements changed:
\stlog{shock_par}

\cgesubsec{Example 2: shocks accumulate}

Shocks build on one another. Two rises of 10 per cent in government purchases
of services raise them by 21 per cent, from 40 to 48.4, as the current
parameters of \manref{cge_parameters} show:
\stlog{shock_accum}
\cmd{cge_shock, reset} removes every shock, those of Example~1 included:
\stlog{shock_reset}

\cgesubsec{Example 3: fixed variables}

Under the default closure labour supply and foreign savings are fixed, and can
be shocked; the exchange rate is free, and cannot:
\stlog{shock_var}
The solution keeps the fixed variables at their new values:
\stlog{shock_solve}

\cgesubsec{Example 4: parameters the model does not have}

The standard synthetic SAM has no health insurance fund, and \cmd{trgr} is a
scalar:
\stlog{shock_need}

\cgesubsec{Example 5: health shocks}

In the health layout of \manref{cge_synthetic}, with a social health insurance
fund mapped by \cmd{cge_map, insurance(shi)}, government purchases of health
services rise by 20 per cent, the contribution rate on labour income by half,
and labour efficiency falls by 2 per cent:
\stlog{shock_health}
The higher contribution rate raises the fund's revenue, and the fund buys more
health services with it.

\cgesec{Methods and formulas}

Let $\theta_e$ be the current value of element $e$ of a parameter, or of a
fixed variable, and $E$ the elements given. \cmd{cge_shock} sets
\begin{equation*}
\theta_e \leftarrow s\,\theta_e \quad\text{(\cmd{scale(}$s$\cmd{)})}, \qquad
\theta_e \leftarrow v \quad\text{(\cmd{value(}$v$\cmd{)})}, \qquad e \in E,
\end{equation*}
except that \cmd{value(}$v$\cmd{)} for \cmd{tf} sets
$\V{tf}_{fa} = v$ for every activity $a$ that uses factor $f$ in the SAM and
leaves the others at zero. \cmd{cge_shock, reset} sets every parameter and
fixed variable to its base value, $\theta = \theta^0$. \manref{cge_solve} holds
the fixed variables at their current values and uses the current parameters;
continuation in $n$ steps solves at
$\theta^0 + (k/n)(\theta - \theta^0)$, $k = 1, \dots, n$.

The parameters enter the equations of \chapref{chap:equations} and of the
extensions: VAT and public debt in \chapref{chap:fiscal}, export demand in
\chapref{chap:trade}, the labour force and wages in \chapref{chap:labour},
health outcomes in \chapref{chap:health}, taxes on factors in
\secref{sec:factortax}, and the rents in \secref{sec:rents}.

\cgesec{Also see}

\begin{alsosee}
\seeentry{cge_closure}{Choose the model closure}
\seeentry{cge_solve}{Solve the model}
\seeentry{cge_constraint}{Limits that bind in some scenarios and not in others}
\seeentry{cge_parameters}{The calibrated parameters as a dataset}
\seechapter{chap:equations}{The equations of the standard model}
\seechapter{chap:closures}{Closures}
\end{alsosee}
